\section{Metodologia}
\label{sec:metodologia}

Esta seção formaliza o protocolo experimental. A
\S{}\ref{sec:protocolo} fixa o protocolo \emph{paper-faithful}
(\emph{split}, validação cruzada, métrica e encoding) compartilhado por
todos os modelos; em seguida descrevemos as extensões de termômetros na
\texttt{wisardpkg} e, por fim, a implementação e o ambiente de
reprodutibilidade (pré-processamento, hardware e artefatos).

\subsection{Protocolo \emph{paper-faithful}}
\label{sec:protocolo}

Para que a comparação entre os \emph{baselines} clássicos e a família
WiSARD seja direta, todos os experimentos seguem o mesmo protocolo,
derivado da Seção VI--A do paper Alsaedi 2020
(Tabela~\ref{tab:protocolo}).

\begin{table}[ht]
  \centering
  \caption{Protocolo \emph{paper-faithful} aplicado a todos os 13
           modelos avaliados (8 \emph{baselines} + 5 WiSARDs).}
  \label{tab:protocolo}
  \small
  \renewcommand{\arraystretch}{1.15}
  \begin{tabular}{p{0.27\linewidth} p{0.32\linewidth} p{0.32\linewidth}}
    \toprule
    \textbf{Item} & \textbf{Paper \S{}VI--A} & \textbf{Nosso} \\
    \midrule
    Split & 80\% \emph{train} / 20\% \emph{test} \emph{held-out} & \texttt{StratifiedShuffleSplit (test\_size=0.20, random\_state=42)} \\
    Validação cruzada & $k=4$ sobre o \emph{train} & \texttt{StratifiedKFold (n\_splits=4, random\_state=42)} \\
    F1 final & Média sobre os 4 \emph{folds} & idem \\
    Métrica F1 & F-score (\S{}VI-B1, não especifica \emph{average}) & \texttt{f1\_weighted} (ancorada no paper + ajuste empírico, \S{}\ref{sec:protocolo_metrica}) \\
    Encoding categórico & \texttt{LabelEncoder} & \texttt{LabelEncoder} cru, \emph{sem} \texttt{clean\_df} (ver \S{}\ref{sec:discussao_clean}) \\
    Temporal & Drop \texttt{date}/\texttt{time}/\texttt{timestamp} & \texttt{drop\_temporal\_leak} \\
    Normalização & \texttt{MinMaxScaler [0, 1]} & idem \\
    Imputação (combined) & Mediana por coluna & idem (apenas numéricas) \\
    \bottomrule
  \end{tabular}
\end{table}

\subsubsection{Por que \emph{F1 weighted}}
\label{sec:protocolo_metrica}

A Seção VI--B1 do paper define F-score com a fórmula binária padrão
sobre TP/FP/FN (precisão e revocação da classe de ataque). A escolha do
\emph{average} do \texttt{sklearn} (\texttt{'binary'}, \texttt{'macro'}
ou \texttt{'weighted'}) muda os valores reportados em mais de 0{,}3
pontos nas bases de maior desbalanceamento, sendo decisiva para qualquer
reprodução. O texto do paper \emph{sugere} a leitura \texttt{weighted}:
na discussão multi-classe da Tabela 13 ele declara calcular ``\emph{a
weighted average\ldots\ multiplying each class with a weight factor
(i.e., the \# instances with that target class)}'', que é exatamente o
\texttt{average='weighted'} do \texttt{sklearn}. Como o paper não publica
a chamada exata para as tarefas binárias (Tabelas 10--11), estendemos
essa leitura e a confirmamos empiricamente: nossa varredura das três
variantes contra essas tabelas mostra \texttt{weighted} como o melhor
ajuste (Tabela~\ref{tab:f1_metric}):

\begin{table}[ht]
  \centering
  \caption{Ajuste empírico das três variantes de F1 às Tabelas 10--11
           (\emph{per-device}) do paper. $\Delta$ = nossa reprodução $-$ paper.}
  \label{tab:f1_metric}
  \small
  \begin{tabular}{lrrrr}
    \toprule
    \textbf{Variante} & \textbf{mean $\Delta$} & \textbf{$|\Delta|<0{,}05$} & \textbf{$|\Delta|<0{,}10$} & \textbf{$|\Delta|<0{,}15$} \\
    \midrule
    \texttt{binary}   & $+0{,}102$ & 38\% & 52\% & 68\% \\
    \texttt{macro}    & $-0{,}032$ & 48\% & 71\% & 84\% \\
    \texttt{weighted} & $\mathbf{-0{,}009}$ & \textbf{48\%} & \textbf{75\%} & \textbf{91\%} \\
    \bottomrule
  \end{tabular}
\end{table}

\texttt{weighted} apresenta o melhor ajuste empírico (média de $\Delta$ mais
próxima de zero e maior fração dentro de tolerância), em concordância
com a definição explícita do paper. Adotamos \texttt{weighted} para
todas as comparações; o \emph{runner} \texttt{run\_baselines\_paper.py}
também grava \texttt{f1} (\emph{average} \texttt{binary}) e \texttt{f1\_macro} para auditoria.

\subsubsection{Por que \emph{sem} \texttt{clean\_df}}

A função \texttt{clean\_df} (apresentada na \S{}\ref{sec:dataset})
neutraliza os \emph{leaks} categóricos identificados. Em uma análise
puramente metodológica, ela seria mandatória. Mas \emph{para
reproduzir o paper}, ela é o passo errado:

\begin{itemize}[leftmargin=*,itemsep=0pt]
  \item Com \texttt{clean\_df}: \texttt{'high'} e \texttt{'high\ '}
        viram a mesma categoria; \texttt{Fridge} e \texttt{Garage\_Door}
        perdem o sinal do encoding e colapsam para predição
        majoritária. Nossos números ficam \emph{muito abaixo} do
        paper.
  \item Sem \texttt{clean\_df}: o \texttt{LabelEncoder} cru preserva
        o encoding (incluindo as variantes \emph{whitespace}); o
        modelo aprende o que aprenderia no dataset original do paper
        (assumindo que o paper de fato encodou da forma descrita em
        sua \S{}VI-A1, ``\emph{high$\rightarrow$1, low$\rightarrow$0}'').
\end{itemize}

A consequência honesta é que nossos modelos podem ter \emph{mais}
sinal do que os do paper em algumas bases (porque os CSVs públicos
têm variantes \emph{whitespace} que o paper provavelmente não viu);
essa diferença aparece nos $\Delta$ positivos reportados em
\S{}\ref{sec:reproducao}.

\subsection{Extensões na \texttt{wisardpkg}}

A biblioteca \texttt{wisardpkg}~\cite{wisardpkg} foi estendida com
termômetros ajustados (\emph{fitted}) para suportar a variedade de
distribuições do TON\_IoT. O \emph{sweep} usa um conjunto curado de
seis variantes \emph{não supervisionadas}, aplicadas uniformemente a
todas as \emph{features}:

\begin{itemize}[leftmargin=*,itemsep=0pt]
  \item \texttt{linear}: cortes equiespaçados (termômetro clássico).
  \item \texttt{distributive}: discretização por quantis empíricos
        da \emph{feature}.
  \item \texttt{gaussian}: centros e largura derivados de
        $(\mu, \sigma)$ estimados no \emph{train}.
  \item \texttt{exponential} e \texttt{logarithmic}: ajustados para
        distribuições de cauda longa (codes Modbus).
  \item \texttt{non-uniform}: \texttt{DynamicThermometer} cujo
        orçamento de \emph{bits} é alocado por \emph{feature} de forma
        proporcional à dispersão (desvio-padrão).
\end{itemize}

Variantes \emph{label-aware} (\texttt{supervised}) e o
\texttt{stochastic} ficaram fora da comparação \emph{unsupervised}: a
primeira usa os rótulos no ajuste, e o segundo (cujo
\texttt{optimize()} não era acionado no \emph{pipeline}) reduzia-se
ao \texttt{distributive}.

\subsection{Implementação e reprodutibilidade}

Todas as funções de pré-processamento vivem em
\texttt{consolidado\_utils.py}:
\texttt{clean\_df} (opcional, não usado na reprodução),
\texttt{drop\_temporal\_leak},
\texttt{build\_combined\_dataset} (concatena as sete bases para a tarefa
\emph{combined}),
\texttt{presence\_bit\_encode},
\texttt{select\_thermometer\_type} e
\texttt{build\_per\_feature\_thermometer} (\emph{dispatch} automático
de termômetros para os modelos WiSARD-family).

Todos os experimentos rodam em um único \emph{venv} Python 3.13
descrito em \texttt{notebooks/toniot/requirements.txt}. As 288
execuções de \emph{baseline} consomem cerca de 1h em
CPU \emph{single-core}; o \emph{sweep} WiSARD-family (\emph{grid}
estendido com tetos de capacidade unificados, \S{}\ref{sec:wisard_fair}),
com 30\,020 configurações e cerca de 150\,mil \emph{fits}, roda em
aproximadamente 6--7h de \emph{wall-clock}, com as sete bases
\emph{per-device} executadas em paralelo. Resultados em
\texttt{results/\_consolidado/} (\emph{baselines} em
\texttt{baselines\_paper.jsonl}; \emph{sweep} em
\texttt{grid/wisard\_grid\_*.jsonl}); o \emph{notebook}
\texttt{06\_consolidado\_grupos.ipynb} é o relatório de referência
re-executável. Todo o código, os \emph{notebooks} e a especificação de
ambiente que produzem estes resultados estão disponíveis em um
repositório público: \url{https://github.com/muanlartins/masters/tree/toniot-wisard-repro}.
